\section{From reduced codegree to the decorated Gram}

The arithmetic codegree theorem counts reduced cells \(c=(s,j)\), whereas a
full physical output also carries \(\gamma\). The \(\gamma\)-multiplicity
must not be silently inserted into or removed from a support count.

For a reduced cell \(c=(s,j)\), let \(\Gamma_c\) be the ambient set of
opposite labels occurring over that cell. For a row \(m\), define the
decorated cell vector, extended by zero on unused labels,
\[
 V_{m,c}:=
 \bigl(b_{(\gamma,s,j),m}\bigr)_\gamma
\]
in \(\ell^2(\Gamma_c)\), with
\[
 E_m(c):=\|V_{m,c}\|^2.
\]
Let \(\mathcal C_{m,n}(R)\) be the common reduced cells of rows \(m,n\)
in the represented block and put
\[
 C_{m,n}^{\rm red}(R):=|\mathcal C_{m,n}(R)|.
\]
This symmetric represented codegree is bounded by the stronger rooted
codegree \(c_{m,n}^{\rm red}(R)\) of
\cref{sec:tpc66-pair-codegree}, in either row ordering.
Then the normalized Gram entry has the exact disintegration
\begin{equation}\label{eq:tpc66-decorated-gram}
 K_X(m,n)
 =
 \frac{1}{\sqrt{T_m(R)T_n(R)}}
 \sum_{c\in\mathcal C_{m,n}(R)}
 \langle V_{m,c},V_{n,c}\rangle.
\end{equation}
Whenever both cell energies are nonzero, define the opposite overlap
\begin{equation}\label{eq:tpc66-opposite-overlap}
 \omega_{m,n}(c)
 :=
 \frac{|\langle V_{m,c},V_{n,c}\rangle|}
 {\sqrt{E_m(c)E_n(c)}}\in[0,1].
\end{equation}
Thus a reduced cell decorated by many \(\gamma\)'s is one Gram correlation,
not automatically many independent phases.

\begin{definition}[Cell diffuseness]
For an active row, set
\begin{equation}\label{eq:tpc66-diffuseness}
 \rho_m(R):=\max_c\frac{E_m(c)}{T_m(R)}\in(0,1].
\end{equation}
\end{definition}

\begin{proposition}[Codegree--diffuseness bridge]
\label{prop:tpc66-codegree-diffuseness}
For distinct active rows,
\begin{equation}\label{eq:tpc66-pair-gram-bound}
 \boxed{\quad
 |K_X(m,n)|
 \le C_{m,n}^{\rm red}(R)
 \sqrt{\rho_m(R)\rho_n(R)}.
 \quad}
\end{equation}
On the critical block \(R\asymp J\), the arithmetic codegree theorem therefore
gives
\begin{equation}\label{eq:tpc66-critical-pair-bound}
 |K_X(m,n)|
 \le X^{o(1)}\sqrt{\rho_m(R)\rho_n(R)}
\end{equation}
uniformly for distinct rows.
\end{proposition}

\begin{proof}
Each summand in \eqref{eq:tpc66-decorated-gram} has magnitude at most
\[
 \sqrt{E_m(c)E_n(c)}
 \le
 \sqrt{\rho_mT_m\,\rho_nT_n}.
\]
There are at most
\(C_{m,n}^{\rm red}(R)\) common reduced cells. Divide by
\(\sqrt{T_mT_n}\).
\end{proof}

\begin{corollary}[Concentration inverse statement]
If \(|K_X(m,n)|\ge\eta>0\), then
\[
 \rho_m(R)\rho_n(R)
 \ge\frac{\eta^2}{(C_{m,n}^{\rm red}(R))^2}.
\]
At \(R\asymp J\), a large off-diagonal correlation therefore forces both rows
to place nonnegligible normalized energy in a small family of common
arithmetic cells.
\end{corollary}

Neither TPC-64 nor TPC-65 proves a fixed-power upper bound for \(\rho_m(R)\).
Support cardinality cannot replace it: a row may have many nonzero cells and
place almost all energy in one of them.

\subsection{Why the residual M\"obius sign is not cancellation}
On the squarefree terminal support,
\[
 \mu(d_m)\mu(r)=\mu(s),\qquad s=d_mr.
\]
The source sign \(\mu(d_m)\) and terminal sign \(\mu(r)\) therefore combine
in the literal coefficient into the output sign \(\mu(s)\), common to every
row contributing to \((\gamma,s,j)\).
Multiplication of one output coordinate by a common unit scalar leaves every
Gram entry unchanged. Thus the residual represented Gram cannot obtain a
parity saving from these M\"obius factors. The fixed-complexity mask analysis
does not furnish small decorated overlap either: its canonical-cell theorem
instead retains \(1-o(1)\) of the relevant product energy across the static
mask, which is a different quantity from
\(\omega_{m,n}(c)\) \citep{WangTPC50}. Genuine signed two-affine
cancellation belongs to the later native/reassembly map studied around
TPC-61 \citep{WangTPC61}, not to the signless residual skeleton isolated
here.
